\documentclass[10pt,a4paper]{article}

\usepackage[T1]{fontenc}
\usepackage{mathpazo}
\usepackage{amsmath}
\usepackage{graphicx}
\usepackage[margin=18mm,top=20mm,bottom=18mm,headsep=7mm]{geometry}
\usepackage{microtype}
\usepackage{fancyhdr}
\usepackage{titlesec}
\usepackage{enumitem}
\usepackage{tabularx}
\usepackage{booktabs}
\usepackage[table]{xcolor}
\usepackage{array}
\usepackage[hidelinks]{hyperref}

\definecolor{ink}{HTML}{172033}
\definecolor{accent}{HTML}{334E68}
\definecolor{pale}{HTML}{E8EEF3}
\definecolor{mid}{HTML}{B8C7D3}
\definecolor{rulegray}{HTML}{778899}

\pagestyle{fancy}
\fancyhf{}
\fancyhead[L]{\small\sffamily\bfseries\color{ink} COL876}
\fancyhead[R]{\small\sffamily\color{ink} Page \thepage}
\renewcommand{\headrulewidth}{0.6pt}
\renewcommand{\headrule}{\hbox to\headwidth{\color{rulegray}\leaders\hrule height \headrulewidth\hfill}}
\setlength{\headheight}{20pt}

\titleformat{\section}
  {\large\bfseries\color{ink}}
  {}{0pt}{}
  [\vspace{-0.35em}\color{mid}\titlerule]
\titlespacing*{\section}{0pt}{0.72em}{0.32em}

\setlist[itemize]{leftmargin=1.25em,itemsep=0.14em,topsep=0.18em,parsep=0pt}
\setlength{\parindent}{0pt}
\setlength{\parskip}{0.3em}
\renewcommand{\arraystretch}{1.1}

\newcommand{\markcell}{\cellcolor{accent}\color{white}\bfseries\centering\arraybackslash $\bullet$}
\newcommand{\softcell}{\cellcolor{pale}}

\begin{document}

% Cover page
\thispagestyle{fancy}
\vspace*{0.045\textheight}

\begin{center}
  {\sffamily\small\bfseries\color{accent}\MakeUppercase{Project Proposal}}\\[0.75cm]

  {\fontsize{38}{43}\selectfont\bfseries\color{ink} KILT}\\[0.28cm]
  {\fontsize{20}{25}\selectfont\bfseries\color{accent} Kripke in Lean Toolkit}\\[0.42cm]
  {\Large\itshape A Verified LTL Model-Checking Toolkit in Lean}\\[0.72cm]

  {\color{mid}\rule{0.82\textwidth}{0.9pt}}\par
  \vspace{0.8cm}

  \begin{minipage}[t]{0.28\textwidth}
    \vspace{0pt}
    \centering
    \includegraphics[width=0.92\linewidth]{lean-logo-official.png}\\[0.75cm]
    {\large\bfseries Yash Rawat}\\[0.18cm]
    {\large 2023CS50334}\\[0.65cm]
    {\normalsize COL876: Special Topics\\[-0.02cm]in Formal Methods}\\[0.18cm]
    {\normalsize September 2026}
  \end{minipage}
  \hfill
  \begin{minipage}[t]{0.60\textwidth}
    \vspace{0.08cm}
    \centering
    \fcolorbox{mid}{pale}{%
      \begin{minipage}{0.92\linewidth}
        \centering\vspace{0.36cm}
        {\small\sffamily\bfseries\color{accent} A LINEAR-TIME REQUIREMENT}\\[0.3cm]
        {\Large $\mathbf{G}(\mathit{request}\rightarrow\mathbf{F}\,\mathit{grant})$}\\[0.28cm]
        {\small\color{ink}Every request is eventually followed by a grant.}\\[0.36cm]
      \end{minipage}%
    }\\[0.25cm]
    {\scriptsize\sffamily\color{accent}An example liveness property for an infinite execution.}
  \end{minipage}
\end{center}

\vfill
\begin{center}
  \begin{minipage}{0.72\textwidth}
    \centering\small\color{accent}
    A Lean development connecting LTL semantics, Büchi automata, finite-state
    verification, and kernel-checked correctness proofs.
  \end{minipage}
\end{center}
\vspace*{0.06\textheight}

\newpage

% Main proposal
\section{Motivation and Background}

Linear Temporal Logic (LTL) describes properties of individual, potentially
infinite executions of a transition system. Operators such as next ($X$),
eventually ($F$), globally ($G$), and until ($U$) express safety and liveness
requirements; for example, $G(request\rightarrow F\,grant)$ states that every
request is eventually granted. LTL is widely used to specify reactive behaviour
and is commonly checked by translating a property into an automaton over infinite
words, then analysing its interaction with the system.

\section{Existing Work and Project Gap}

LeanLTL is already a public Lean 4 library and a peer-reviewed ITP 2025 paper. It
supports LTL reasoning over finite and infinite traces, offers a syntax macro,
and provides automation for proofs about temporal formulas [1]. Separately,
CSLib includes automata-theory foundations for finite and infinite words,
including Büchi automata, products, and omega-language results [2,3]. A second
Lean 4 project, LeanearTemporalLogic, formalizes LTL syntax, semantics, and
transition systems; its README lists checked LTL model checking and Büchi
formalization as ambitious future goals and describes the project as work in
progress [4]. Veil is another relevant Lean-native verification framework for
transition systems; I will study its workflow and examples as groundwork, while
treating it as a reference rather than a required dependency [5].

These projects establish substantial foundations, so the proposal does not
claim that LTL or Büchi automata are new to Lean, or that no Lean LTL checker is
under development. The intended contribution is narrower: a CSLib-native,
reusable path from a bounded LTL fragment through Büchi acceptance and product
emptiness to a verified model-checking result. I will use LeanLTL as prior work
and a semantic comparison point, rather than simply wrapping it as the project
implementation; a small formula layer will be introduced only as needed to
drive the automata construction and CSLib APIs. Where feasible, I will prove
that its satisfaction relation agrees with LeanLTL's trace semantics. During
the initial audit I will inspect LeanearTemporalLogic's current implementation
and discuss overlap through Lean Zulip, then narrow the module-level contribution
to a clearly distinct and useful CSLib extension.

\section{Goals and Expected Deliverables}

\begin{itemize}
  \item Audit LeanLTL, CSLib's automata APIs, LeanearTemporalLogic, and Veil;
        consult the Lean community about overlap, then select a clearly bounded
        LTL fragment and design only the formula interface needed for CSLib
        translation. Prove a compatibility result with LeanLTL's trace
        semantics where feasible.
  \item Develop an LTL-to-Büchi (or generalized Büchi) construction for the
        selected fragment, reusing CSLib's automata interfaces where appropriate.
  \item Build the product of a finite Kripke structure with the property
        automaton and an executable emptiness check for accepting infinite runs.
  \item Prove the central correspondence between the checker and the declared
        trace semantics. Correctness of the semantic and automata-theoretic core
        will take priority over extensions.
  \item If the core is complete, add a lightweight Lean metaprogramming frontend
        (notation, macro, or command) to make properties and checks convenient to
        write. Proof-producing automation and an initial CTL* study are further
        directions, not required deliverables.
  \item Deliver organized Lean source, examples/tests (including representative
        safety and liveness properties), and a short write-up of the design,
        proofs, trust boundary, and limitations.
\end{itemize}

\textbf{Division of labour.} This is an individual project. I will undertake
the formal design, implementation, correctness proofs, interface development,
evaluation, and documentation.

\newpage
\section{Expected Milestones: 24 September--15 November}

The first week is reserved for research and setup. The dated blocks below show
the intended progression; the core semantic correspondence and verified
model-checking result take priority over frontend polish, witness presentation,
and any CTL* extension.

\footnotesize
\setlength{\tabcolsep}{2.7pt}
\begin{tabularx}{\textwidth}{>{\raggedright\arraybackslash}X *{7}{>{\centering\arraybackslash}p{1.12cm}}}
  \toprule
  \textbf{Workstream} & \textbf{24--30 Sep} & \textbf{1--7 Oct} &
  \textbf{8--14 Oct} & \textbf{15--21 Oct} & \textbf{22--31 Oct} &
  \textbf{1--7 Nov} & \textbf{8--15 Nov} \\
  \midrule
  Literature/API audit, Lean setup, scoped design
    & \markcell & & & & & & \\
  LTL syntax/semantics and Kripke execution model
    & & \markcell & \softcell & & & & \\
  CSLib Büchi interface and formula-to-automaton core
    & & \softcell & \markcell & \softcell & & & \\
  Product construction and accepting-run emptiness
    & & & \softcell & \markcell & \softcell & & \\
  Correctness theorem and verified examples
    & & & & \softcell & \markcell & \softcell & \\
  Lean frontend, evaluation, and write-up
    & & & & & & \softcell & \markcell \\
  Optional CTL* feasibility study / extensions
    & & & & & & & \softcell \\
  \bottomrule
\end{tabularx}

\vspace{0.5em}
\small
The exact supported fragment and reuse boundary will be fixed after the
September audit. Witnesses and proof automation will be attempted only after the
verified checker is working; CTL* is an exploratory extension if the core goals
are complete.

\section{Selected References}
\begin{enumerate}[leftmargin=1.5em,itemsep=0.18em,topsep=0.2em]
  \item E. Vin, K. A. Miller, and D. J. Fremont, ``LeanLTL: A Unifying Framework
        for Linear Temporal Logics in Lean,'' \emph{ITP 2025}, LIPIcs 352,
        Article 37. \href{https://doi.org/10.4230/LIPIcs.ITP.2025.37}{doi:10.4230/LIPIcs.ITP.2025.37};
        code: \url{https://github.com/UCSCFormalMethods/LeanLTL}.
  \item Leanprover, \emph{CSLib: The Lean Computer Science Library}, automata
        modules and contribution guide. \url{https://github.com/leanprover/cslib}.
  \item C.-T. Chou, \emph{Automata Theory in Lean}, results being ported to CSLib.
        \url{https://github.com/ctchou/AutomataTheory}.
  \item M. Pawagi, \emph{LeanearTemporalLogic}, Lean 4 LTL formalization (work in
        progress). \url{https://github.com/mrigankpawagi/LeanearTemporalLogic}.
  \item Verse Lab, \emph{Veil: A Framework for Automated and Interactive
        Verification of Transition Systems}. \url{https://github.com/verse-lab/veil}.
\end{enumerate}

\end{document}
